\section{06.11.2025 -- Matching Day und Selbstpräsentation}

\subsection{Aufgabe}

Beim Treffen simulieren wir eine Selbstpräsentation auf einem Matching Day. Vorbereitung:
\begin{itemize}
    \item Unternehmen auswählen.
    \item Unternehmenstyp notieren: Start-up, kleines Unternehmen oder Konzern.
    \item Branche notieren: Autoindustrie, Pharma, Tech, Finanzen, Bildung, Kunst usw.
    \item Zwei- bis dreiminütige Selbstpräsentation vorbereiten.
\end{itemize}

\subsection{Selbstpräsentation -- CERN}

Heute möchte ich über CERN sprechen. CERN ist ein Forschungszentrum, in dem Physikerinnen und Physiker aus der ganzen Welt viele wissenschaftliche Experimente durchführen.

Hallo zusammen, ich bin Vyron und ich komme aus Griechenland. Ich freue mich sehr, dass ich hier sein kann. Ich interessiere mich für die PhD-Stelle ``Particle Physics und KI'' am CERN, weil Teilchenphysik in den letzten Jahren meine Leidenschaft geworden ist und ich seit letztem Jahr auch eine starke Verbindung zur KI habe.

Ich habe meinen Bachelor in Physik an der Universität Athen gemacht. Dort habe ich meine Thesis am CERN durchgeführt. Genauer gesagt haben wir eine Datenanalyse für seltene Teilchen durchgeführt. Außerdem arbeite ich seit April als Werkstudent beim DLR im Bereich Softwareentwicklung. Die Software wird für Machine-Learning-Analysen verwendet.

Parallel dazu schreibe ich meine Masterarbeit am Max-Planck-Institut für Quantenoptik. Dort entwickle ich Machine-Learning-Algorithmen und versuche, möglichst viele Informationen aus einem medizinischen Datensatz zu gewinnen.

Durch meine Erfahrungen kann man sehen, dass ich alle nötigen Fähigkeiten habe, um diese PhD-Stelle anzustreben. Außerdem bin ich sehr motiviert. Ich freue mich darauf, mein Wissen weiterzuentwickeln. Ich bin überzeugt, dass ich mit meiner Leidenschaft für Physik und KI einen wertvollen Beitrag zu Ihrem Team leisten kann.

\subsection{Matching Day vs. Karrieremesse}

\begin{itemize}
    \item \textbf{Matching Day}: mit mehr Vorbereitung.
    \item \textbf{Karrieremesse}: auch ohne große Vorbereitung möglich.
\end{itemize}

\subsection{Korrekturen und Formulierungen}

\begin{itemize}
    \item \wrong{Auf die andere Seite} $\rightarrow$ \correct{Auf der anderen Seite}
    \item \correct{Auf der einen Seite \dots{} auf der anderen Seite \dots}
    \item \correct{Meiner Meinung nach \dots}
    \item \correct{Ich bin der Meinung, dass \dots}
    \item \correct{Ich finde / denke / glaube, dass \dots}
\end{itemize}

\subsection{Forum 2.9 -- Welche Arbeit passt zu mir?}

\subsubsection{Interessante Wörter}

\begin{vocabtable}
\deword{der Personalchef}{HR manager, head of personnel}{Der Personalchef führt das Bewerbungsgespräch.}
\end{vocabtable}

\subsubsection{Verbesserungen}

\begin{itemize}
    \item \wrong{eine anpassende Veranstaltung} $\rightarrow$ \correct{eine passende Veranstaltung}
    \item \correct{ein geeignetes Event} / \correct{ein geeignetes Ereignis}
    \item \wrong{sich über die Firmen, die daran Interesse haben, gut zu informieren} $\rightarrow$ \correct{sich über die Firmen, an denen man Interesse hat, gut zu informieren}
    \item \correct{sich über die Firmen, für die man sich interessiert, gut zu informieren}
    \item \wrong{eine Herausforderung für Ausländern} $\rightarrow$ \correct{eine Herausforderung für Ausländer}
    \item \wrong{die Mehrheit der deutschen Firmen bevorzugen zu auf Deutsch sprechen} $\rightarrow$ \correct{die Mehrheit der deutschen Firmen bevorzugt deutschsprachige Mitarbeitende}
    \item \correct{die Mehrheit der deutschen Firmen bevorzugt es, auf Deutsch zu sprechen}
    \item \wrong{ein online Jobbörse} $\rightarrow$ \correct{eine Online-Jobbörse}
    \item \wrong{aus der Menge auszustehen} $\rightarrow$ \correct{aus der Menge herauszustechen}
    \item \wrong{eine große Möglichkeit} $\rightarrow$ \correct{eine gute Gelegenheit / eine gute Chance}
    \item \correct{In der Theorie halte ich einen Matching Day für eine gute Gelegenheit, verschiedene Firmen kennenzulernen.}
    \item \wrong{ein Gefühl für die Unternehmenskultur zu finden} $\rightarrow$ \correct{ein Gefühl für die Unternehmenskultur zu bekommen / zu entwickeln}
    \item \wrong{Ich glaube, dass ein Matching Day ist gut} $\rightarrow$ \correct{Ich glaube, dass ein Matching Day gut ist}
\end{itemize}

\subsection{Eigene Formulierungen}

Ich stimme zu, weil Menschen Menschen mögen. Wenn man jemanden persönlich trifft, mag die andere Person einen oft mehr, als wenn sie nur eine Bewerbung liest.
